\documentclass[11pt,letterpaper]{article}
\usepackage[margin=0.82in]{geometry}
\usepackage[T1]{fontenc}
\usepackage[utf8]{inputenc}
\usepackage{lmodern}
\usepackage{microtype}
\usepackage{xcolor}
\usepackage{booktabs}
\usepackage{array}
\usepackage{longtable}
\usepackage{listings}
\usepackage{amsmath}
\usepackage{enumitem}
\usepackage{fancyhdr}
\usepackage{hyperref}
\usepackage{titlesec}

\definecolor{codebg}{RGB}{246,248,250}
\definecolor{codeframe}{RGB}{216,222,228}
\definecolor{accent}{RGB}{28,78,121}
\hypersetup{colorlinks=true,linkcolor=accent,urlcolor=accent}
\lstdefinestyle{pythonstyle}{
  language=Python,
  basicstyle=\ttfamily\small,
  keywordstyle=\color{blue!55!black},
  commentstyle=\color{green!35!black},
  stringstyle=\color{red!45!black},
  backgroundcolor=\color{codebg},
  frame=single,
  rulecolor=\color{codeframe},
  framerule=0.35pt,
  breaklines=true,
  keepspaces=true,
  showstringspaces=false,
  columns=fullflexible,
  xleftmargin=4pt,
  xrightmargin=4pt,
  aboveskip=8pt,
  belowskip=8pt
}
\lstset{style=pythonstyle}
\setlist[itemize]{leftmargin=1.4em,itemsep=0.2em,topsep=0.35em}
\setlist[enumerate]{leftmargin=1.55em,itemsep=0.25em,topsep=0.35em}
\titleformat{\section}{\Large\bfseries\color{black}}{\thesection}{0.65em}{}
\titleformat{\subsection}{\large\bfseries\color{black}}{\thesubsection}{0.65em}{}
\pagestyle{fancy}
\fancyhf{}
\fancyhead[L]{\small Clean and Prepare Data with Python}
\fancyhead[R]{\small Study Guide}
\fancyfoot[C]{\thepage}
\renewcommand{\headrulewidth}{0.3pt}
\setlength{\parindent}{0pt}
\setlength{\parskip}{0.55em}

\title{Clean and Prepare Data for Analysis with Python\\[0.35em]
\large Study Guide for Scientific and Time Series Data}
\author{}
\date{}

\begin{document}
\maketitle
\vspace{-2em}

This guide explains how to turn imperfect observations into an analysis-ready
dataset with NumPy and pandas. The accompanying Python script runs a complete
workflow, while the Jupyter notebook provides guided practice and solutions.
The central example is a light curve with slow drift, transit-like dips,
isolated spikes, uneven sampling, missing values, and measurement noise.

\textbf{Estimated study time:} 60--90 minutes. Use the Python environment at
\nolinkurl{/Users/victorbui/venvs/ai312}. Keep the script and notebook in the same
folder. Run the script with:

\begin{lstlisting}
# Execute the complete course and create the output tables and figures.
/Users/victorbui/venvs/ai312/bin/python \
    Python_Clean_Prepare_Data_Course.py
\end{lstlisting}

\tableofcontents

\section{Learning plan}

The workflow follows the order in which decisions should normally be made.
Inspect the raw records before applying a method. Standardize missing-value
codes before calculating statistics. Detect unusual observations in context.
Interpolate only gaps that the data can support. Apply scaling, detrending, and
smoothing with a stated analytical purpose. Finish with explicit checks and an
audit table.

\begin{center}
\begin{tabular}{@{}p{0.19\linewidth}p{0.33\linewidth}p{0.38\linewidth}@{}}
\toprule
Stage & Main question & Evidence to retain \\
\midrule
Inspection & What does each row and column mean? & Shape, types, units, ranges, time order \\
Missing data & Which markers mean unavailable? & Counts, locations, gap lengths, chosen policy \\
Outliers & Which values violate the local pattern? & Method, window, threshold, flagged rows \\
Interpolation & Which gaps support estimation? & Sample locations, method, maximum gap \\
Normalization & Which comparison needs a common scale? & Fitted mean, spread, range, units \\
Detrending & Which baseline obscures the event? & Model degree, excluded events, residuals \\
Smoothing & Which variation is noise? & Window width, direction, statistic \\
Validation & Did the workflow preserve meaning? & Assertions, row counts, output schema \\
\bottomrule
\end{tabular}
\end{center}

\textbf{Key distinction.} Missing values indicate unavailable measurements.
Outliers are observed values that appear inconsistent with an expected pattern.
Noise is short-scale variation. A trend is structured change over a longer
scale. These conditions require different evidence and should not share one
automatic treatment.

\section{Inspect data before cleaning}

Start by confirming what the table represents. For time series data, verify
that timestamps are ordered and determine whether the spacing is regular. Read
units and instrument documentation. Plot the data because summaries can hide
spikes, gaps, and drift.

\begin{lstlisting}
# Inspect structure, missingness, order, and numeric ranges.
print(data.shape)
print(data.dtypes)
print(data.isna().sum())
print(data.describe())
assert data["time_days"].is_monotonic_increasing

ax = data.plot(x="time_days", y="brightness", figsize=(10, 4))
ax.set(xlabel="Time (days)", ylabel="Relative brightness")
\end{lstlisting}

A sentinel such as \texttt{-999} may represent a missing measurement. Until it
is replaced, the minimum, mean, standard deviation, trend fit, and plots treat
it as a valid number. Only replace a code when metadata or domain knowledge
defines its meaning.

\begin{lstlisting}
# Convert documented codes to the standard numeric missing value.
data["brightness"] = data["brightness"].replace(-999, np.nan)
\end{lstlisting}

\subsection{A useful inspection checklist}

\begin{itemize}
\item Confirm the unit and meaning of every variable used in the analysis.
\item Check duplicate rows and duplicate timestamps separately.
\item Count explicit missing values and documented sentinel codes.
\item Look for impossible values using domain limits.
\item Inspect time order, sample spacing, and gap length.
\item Compare a plot with numeric summaries before changing the data.
\end{itemize}

\section{Detect and handle outliers}

An outlier is unusual relative to a specified reference. A global mean and
standard deviation may work for stable, symmetric data, but they are sensitive
to extreme observations. Quartile rules are more robust. A local signal often
needs a moving reference because a slow trend changes its center over time.

The course uses a rolling median and median absolute deviation (MAD). Within a
local window, let $m$ be the median. The raw MAD is

\[
  \operatorname{MAD} = \operatorname{median}(|x_i-m|).
\]

For nearly normal noise, $1.4826\operatorname{MAD}$ estimates the standard
deviation. A point is flagged when its absolute deviation exceeds a chosen
multiple of that robust scale.

\begin{lstlisting}
# Estimate a local center and robust local scale.
center = values.rolling(41, center=True, min_periods=9).median()
deviation = (values - center).abs()
mad = deviation.rolling(41, center=True, min_periods=9).median()
robust_scale = 1.4826 * mad
outlier = (deviation > 5 * robust_scale) & values.notna()

# Preserve the timestamp and mark the suspect measurement as missing.
cleaned = data.copy()
cleaned.loc[outlier, "brightness"] = np.nan
\end{lstlisting}

The window should cover enough ordinary variation to estimate a stable local
center without spanning unrelated behavior. The threshold controls sensitivity:
a smaller number flags more observations. Review flagged points on the raw plot.
A scientifically meaningful event can be rare and extreme without being an
error.

\textbf{Removal versus marking.} Dropping an outlier row also removes its time
coordinate. Marking its measurement as missing retains the sampling grid and
records where information became unavailable. This often makes later gap logic
clearer.

\section{Work with missing data}

There are three broad strategies. Ignoring missing values during a reduction
retains the measured observations. Removing rows creates a complete subset but
can discard useful values from other columns. Replacing values preserves shape
and alignment but adds estimates that were not observed.

\begin{longtable}{@{}p{0.18\linewidth}p{0.33\linewidth}p{0.39\linewidth}@{}}
\toprule
Strategy & Appropriate use & Main limitation \\
\midrule
\endhead
Ignore in a calculation & A reduction can use available values and the missing mechanism does not invalidate it & Different calculations may handle missing values differently \\
Remove rows & Missing rows are few and removal does not distort the sample or break alignment & Valid values in other columns may be lost \\
Replace with a constant or statistic & The replacement has a documented meaning, or a simple baseline is explicitly required & Variability and relationships can be distorted \\
Interpolate & Nearby observations support an estimate within a short gap & Filled values can look more certain than they are \\
Model-based imputation & A justified model uses other information and uncertainty is tracked & Results depend on model assumptions \\
\bottomrule
\end{longtable}

\begin{lstlisting}
# Locate, count, ignore, remove, or replace missing values.
missing = data["brightness"].isna()
count = int(missing.sum())
mean_available = data["brightness"].mean()
complete_rows = data.dropna(subset=["brightness"])
median_filled = data["brightness"].fillna(data["brightness"].median())
\end{lstlisting}

When a plotted line contains \texttt{NaN}, it breaks at the gap. Dropping those
rows can connect the two measured segments and visually imply continuity. Keep
the original timestamps and a missingness indicator when the gaps matter.

\section{Interpolate missing values}

Interpolation estimates values between known observations. Nearest-neighbor
interpolation produces steps. Linear interpolation connects neighboring values
with straight lines. Cubic splines are smooth but can oscillate between points.
Shape-preserving cubic interpolation limits those oscillations when monotonicity
matters.

For points $(x_1,y_1)$ and $(x_2,y_2)$, linear interpolation at $x$ is

\[
  y(x) = y_1 + \frac{x-x_1}{x_2-x_1}(y_2-y_1).
\]

The sample locations matter. If $x$ is unevenly spaced, interpolating by row
position gives a different answer from interpolating by the actual $x$ values.

\begin{lstlisting}
# Interpolate using the numeric index as the actual sample location.
timed = pd.Series(y.to_numpy(), index=x)
linear_by_x = timed.interpolate(method="index")
\end{lstlisting}

Limit interpolation to gaps supported by neighboring information. The following
pattern fills no more than five consecutive internal missing rows. Boundary
gaps and longer runs remain missing.

\begin{lstlisting}
# Preserve long gaps instead of drawing unsupported straight bridges.
timed = data.set_index("time_days")["brightness"]
missing = timed.isna()
run_id = missing.ne(missing.shift()).cumsum()
run_size = missing.groupby(run_id).transform("sum")
fillable = missing & (run_size <= 5)

candidate = timed.interpolate(method="index", limit_area="inside")
timed.loc[fillable] = candidate.loc[fillable]
\end{lstlisting}

The row limit is an operational rule, not a universal scientific threshold. If
sample spacing varies greatly, measure gap duration as well as row count. Store
an indicator for interpolated rows when later users need to distinguish estimates
from observations.

\section{Shift and scale data}

Normalization makes variables comparable or converts a physical measurement to
a form required by a model. Each transformation answers a different question.
Missing values should normally remain missing during the calculation.

\subsection{Centering}

Centering subtracts the mean:

\[
  x_{c,i} = x_i - \bar{x}.
\]

The centered data have mean zero while retaining their unit and spread.

\subsection{Z-scores}

A z-score measures distance from the mean in standard deviations:

\[
  z_i = \frac{x_i-\bar{x}}{s}.
\]

Specify whether $s$ is the population standard deviation
(\texttt{ddof=0}) or sample standard deviation (\texttt{ddof=1}). This course
uses \texttt{ddof=0} and verifies mean 0 and standard deviation 1.

\subsection{Relative and range scaling}

Relative change can place a measured signal on a dimensionless scale:

\[
  r_i = \frac{x_i-\bar{x}}{\bar{x}}.
\]

Min-max scaling maps observations to $[a,b]$:

\[
  x'_i = a + (b-a)\frac{x_i-x_{\min}}{x_{\max}-x_{\min}}.
\]

\begin{lstlisting}
# Calculate several forms without overwriting the measured values.
mean_value = x.mean()
population_sd = x.std(ddof=0)
centered = x - mean_value
z_score = (x - mean_value) / population_sd
relative = (x - mean_value) / mean_value
scaled_01 = (x - x.min()) / (x.max() - x.min())
\end{lstlisting}

Constant data have zero spread and cannot be divided by their standard deviation
or range. In prediction workflows, calculate normalization parameters on the
training partition, retain them, and apply the same values to validation and
future observations.

\section{Remove unwanted trends}

Detrending separates a slowly changing baseline from the variation under study.
A polynomial baseline of degree $n$ can be fitted with \texttt{np.polyfit} and
evaluated with \texttt{np.polyval}. Subtracting the fitted baseline produces the
residual signal.

\begin{lstlisting}
# Exclude missing values and known events from the baseline fit.
valid = values.notna() & baseline_mask
time_center = time[valid].mean()
time_scale = time[valid].std(ddof=0)
scaled_time = (time - time_center) / time_scale

coefficients = np.polyfit(scaled_time[valid], values[valid], 3)
trend = np.polyval(coefficients, scaled_time)
detrended = values - trend
\end{lstlisting}

Centering and scaling time improves numerical conditioning, particularly for
higher-degree polynomials or large timestamp values. Exclude known events from
the baseline fit when the events should survive in the residual. Otherwise, the
baseline can bend toward them and reduce their apparent depth.

Model degree creates a bias-flexibility tradeoff. A line may leave curvature. A
high-degree polynomial may follow noise or absorb the event. Compare residual
structure, scientific plausibility, and performance on held-out baseline regions.
The smallest training residual is insufficient evidence by itself.

\section{Smooth noisy data}

Smoothing reduces short-scale variation so broader patterns become easier to
see. It also changes the signal. Record the method, window size, direction, and
boundary behavior.

\begin{center}
\begin{tabular}{@{}p{0.25\linewidth}p{0.28\linewidth}p{0.37\linewidth}@{}}
\toprule
Method & Strength & Limitation \\
\midrule
Moving mean & Reduces zero-mean random noise & Sensitive to spikes and blurs sharp changes \\
Moving median & Resists isolated extreme values & Can flatten small-scale continuous variation \\
Gaussian-weighted mean & Gives nearby values greater weight & Requires a bandwidth choice \\
Local regression & Follows smooth nonlinear patterns & More computation and tuning \\
Polynomial smoothing filter & Can preserve peaks better than a simple mean & Window and polynomial order must be compatible \\
\bottomrule
\end{tabular}
\end{center}

\begin{lstlisting}
# Sample-count windows for regularly spaced observations.
centered_mean = series.rolling(5, center=True, min_periods=1).mean()
trailing_median = series.rolling(4, min_periods=1).median()

# Reverse twice to calculate a leading four-point median.
leading_median = (
    series.iloc[::-1].rolling(4, min_periods=1).median().iloc[::-1]
)
\end{lstlisting}

A centered five-point window uses two previous values, the current value, and
two future values. A trailing four-point window uses the current value and three
previous values. Trailing windows support live monitoring because they do not
use future data. Leading and centered windows suit retrospective analysis.

For irregular sampling, a fixed number of rows does not represent a fixed time
span. Use a time-aware window. The course script includes
\texttt{time\_window\_smooth}, which defines separate durations before and after
each observation.

Larger windows suppress more noise, but they also reduce peak height, broaden
transitions, and can erase short events. Compare several widths against the
known time scale of the phenomenon.

\section{Build a reviewable pipeline}

A reliable pipeline retains raw inputs, creates named intermediate columns, and
tests expected invariants. Avoid overwriting the only copy of a measurement.
Export an audit table alongside the prepared observations.

\begin{lstlisting}
# Check properties that must remain true after preprocessing.
assert len(final_data) == len(raw)
assert final_data["time_days"].is_monotonic_increasing
assert not (final_data["brightness_clean"] == -999).any()
assert final_data.loc[long_gap, "brightness_clean"].isna().all()
assert abs(z_score.mean()) < 1e-10
assert abs(z_score.std(ddof=0) - 1) < 1e-10

# Export without turning the DataFrame index into an extra data column.
final_data.to_csv("output/prepared_light_curve.csv", index=False)
audit.to_csv("output/quality_audit.csv", index=False)
\end{lstlisting}

\subsection{Suggested output schema}

\begin{center}
\begin{tabular}{@{}p{0.28\linewidth}p{0.62\linewidth}@{}}
\toprule
Column & Meaning \\
\midrule
\texttt{time\_days} & Original sample location in days \\
\texttt{brightness\_clean} & Sentinel and outlier handling applied; long gaps retained \\
\texttt{relative\_brightness} & Centered brightness divided by its mean \\
\texttt{trend} & Cubic baseline estimated outside event intervals \\
\texttt{detrended} & Relative brightness minus the fitted trend \\
\texttt{smoothed} & Centered mean over a 0.10-day window \\
\texttt{trailing\_median} & Causal median over the current and prior 0.10 days \\
\bottomrule
\end{tabular}
\end{center}

\section{Practice questions}

\begin{enumerate}
\item Why must a sentinel code be standardized before calculating a mean or
fitting a trend?
\item When does a rolling outlier detector provide a better reference than a
global threshold?
\item Why can dropping missing rows make a time-series line plot misleading?
\item For irregular sample locations, why does linear interpolation by row
position differ from interpolation by the measured location?
\item Give one reason to retain a long gap instead of interpolating it.
\item Which normalization keeps the original unit? Which transformations are
dimensionless?
\item Why should known transit-like events be excluded from a baseline fit?
\item Which smoothing window is appropriate for live monitoring, and why?
\item Why can a 17-point moving mean obscure information that a 3-point moving
mean preserves?
\item Name four assertions that should run before exporting the prepared data.
\end{enumerate}

\subsection*{Answer guide}

\begin{enumerate}
\item The code otherwise acts as a real measurement and can dominate summaries
and fitted models.
\item It helps when the signal center changes over time and unusual values must
be judged against nearby observations.
\item Removing the missing row removes the gap marker, so the plot can connect
separate measured segments.
\item Row position assumes equal spacing; actual-location interpolation weights
the estimate by the real distances between neighboring samples.
\item A long gap may contain unobserved structure, and a straight bridge can be
mistaken for measured continuity.
\item Centering retains the unit. Z-scores, relative change, and unit-interval
scaling are dimensionless.
\item Including the events can pull the baseline toward them and reduce their
apparent depth in the residual.
\item A trailing window, because it uses only current and earlier observations.
\item The wider window averages across a longer interval and therefore blurs or
flattens short features.
\item Examples include row count, time order, sentinel removal, long-gap
preservation, expected schema, and normalization properties.
\end{enumerate}

\section{Quick reference}

\small
\begin{longtable}{@{}p{0.27\linewidth}p{0.64\linewidth}@{}}
\toprule
Goal & Python pattern \\
\midrule
\endhead
Locate missing values & \texttt{series.isna()} \\
Count missing values & \texttt{series.isna().sum()} \\
Standardize a sentinel & \texttt{series.replace(-999, np.nan)} \\
Remove incomplete rows & \texttt{df.dropna(subset=["value"])} \\
Fill with a median & \texttt{series.fillna(series.median())} \\
Interpolate by location & \texttt{pd.Series(y, index=x).interpolate(method="index")} \\
Centered moving mean & \texttt{series.rolling(5, center=True).mean()} \\
Trailing moving median & \texttt{series.rolling(4).median()} \\
Center values & \texttt{x - x.mean()} \\
Calculate z-scores & \texttt{(x - x.mean()) / x.std(ddof=0)} \\
Scale to $[0,1]$ & \texttt{(x - x.min()) / (x.max() - x.min())} \\
Fit a polynomial & \texttt{coeff = np.polyfit(x, y, degree)} \\
Evaluate a polynomial & \texttt{trend = np.polyval(coeff, x)} \\
Select valid rows & \texttt{valid = series.notna()} \\
Preserve a copy & \texttt{cleaned = raw.copy()} \\
Validate a requirement & \texttt{assert condition} \\
Export a CSV & \texttt{df.to\_csv("prepared.csv", index=False)} \\
\bottomrule
\end{longtable}
\normalsize

\section*{Completion criteria}

You have completed the material when you can run both course formats, explain
why each row was removed or estimated, preserve the intended long gap, compare
at least two normalization and smoothing choices, and export a dataset whose
quality checks pass. A preprocessing method is complete only when its assumptions
and effects are visible to the next analyst.

\end{document}
